\documentclass[hidelinks,12pt]{article}
\usepackage[left=0.25cm,top=1cm,right=0.25cm,bottom=1cm]{geometry}
%\usepackage[landscape]{geometry}
\textwidth = 20cm
\hoffset = -1cm
\usepackage[utf8]{inputenc}
\usepackage[spanish,es-tabla, es-lcroman]{babel}
\usepackage[autostyle,spanish=mexican]{csquotes}
\usepackage[tbtags]{amsmath}
\usepackage{nccmath}
\usepackage{amsthm}
\usepackage{amssymb}
\usepackage{mathrsfs}
\usepackage{graphicx}
\usepackage{subfig}
\usepackage{caption}
%\usepackage{subcaption}
\usepackage{standalone}
\graphicspath{{Imagenes/}{../Imagenes/}}
\usepackage[outdir=./Imagenes/]{epstopdf}
\usepackage{siunitx}
\usepackage{physics}
\AtBeginDocument{\RenewCommandCopy\qty\SI}
\ExplSyntaxOn
\msg_redirect_name:nnn { siunitx } { physics-pkg } { none }
\ExplSyntaxOff
\usepackage{color}
\usepackage{float}
\usepackage{hyperref}
\usepackage{multicol}
\usepackage{multirow}
%\usepackage{milista}
\usepackage{anyfontsize}
\usepackage{anysize}
%\usepackage{enumerate}
\usepackage[shortlabels]{enumitem}
\usepackage{capt-of}
\usepackage{bm}
\usepackage{mdframed}
\usepackage{relsize}
\usepackage{placeins}
\usepackage{empheq}
\usepackage{cancel}
\usepackage{pdfpages}
\usepackage{wrapfig}
\usepackage[flushleft]{threeparttable}
\usepackage{makecell}
\usepackage{fancyhdr}
\usepackage{tikz}
\usepackage{bigints}
\usepackage{tcolorbox}
\tcbuselibrary{breakable}
\usepackage{scalerel}
\usepackage{pgfplots}
\usepackage{pdflscape}
\usepackage{enumitem}
\pgfplotsset{compat=1.16}
\spanishdecimal{.}
\renewcommand{\baselinestretch}{1.5}
\def\scaleint#1{\vcenter{\hbox{\scaleto[3ex]{\displaystyle\int}{#1}}}}
\def\scaleoint#1{\vcenter{\hbox{\scaleto[3ex]{\displaystyle\oint}{#1}}}}
\def\scaleiint#1{\vcenter{\hbox{\scaleto[3ex]{\displaystyle\iint}{#1}}}}
\def\scaleiiint#1{\vcenter{\hbox{\scaleto[3ex]{\displaystyle\iiint}{#1}}}}
\def\bs{\mkern-12mu}

\newcommand{\Cancel}[2][black]{{\color{#1}\cancel{\color{black}#2}}}

% \newcommand{\qed}{\tag*{$\blacksquare$}}
\renewcommand{\qed}{\hfill\blacksquare}

\newcommand{\pderivada}[1]{\ensuremath{{#1}^{\prime}}}
\newcommand{\sderivada}[1]{\ensuremath{{#1}^{\prime \prime}}}
\newcommand{\tderivada}[1]{\ensuremath{{#1}^{\prime \prime \prime}}}
\newcommand{\nderivada}[2]{\ensuremath{{#1}^{(#2)}}}

\title{Lista de ejercicios del Tema 2 \\[0.3em]  \large{Matemáticas Avanzadas de la Física}\vspace{-3ex}}
\author{M. en C. Gustavo Contreras Mayén}
\date{ }

\begin{document}
\vspace{-4cm}
\maketitle

\fontsize{14}{14}\selectfont

\textbf{Indicaciones: } Se te pide gentilmente que resuelvas de manera detallada, clara y ordenada los siguientes ejercicios, el puntaje que otorga cada enunciado es de \textbf{1 punto}. En caso de que requieras apoyarte en alguna propiedad, si fue vista en clase, solo indícalo, pero si esa propiedad aunque esté relacionada al ejercicio y no se haya mencionado en clase, habrá que demostrarla debidamente.

\subsection*{Separación de variables.}

\begin{enumerate}
    \item Encuentra la solución a la ecuación de onda:
        \begin{align*}
            u_{tt} = 2 \, u_{xx} 
        \end{align*}
    en $[0, 2 \, \pi]$ que satisface las CDF: $u (0, t) = 0$ y $u_{x} (2 \, \pi, t) = 0$
    \item Encuentra la solución a la ecuación de Laplace:
        \begin{align*}
            u_{xx} + u_{yy} = 0
        \end{align*}
    en un cuadrado unitario que satisface las CDF:
        \begin{align*}
            u (0, y) &= 0 \hspace{1cm} u (1, y) = g (y) \\[0.5em]
            u (x, 0) &= 0  \hspace{1cm} u (x, 1) = 0
        \end{align*}
        \item Clasifica las siguientes ecuaciones ya sea como hiperbólicas, parabólicas o elípticas:
            \begin{enumerate}
            \item $u_{yy} + u_{xy} + u_{xx} = 0$
            \item $3 \, u_{xx} + 2 \, u_{xy} + 5 \, u_{yy} = 0$
            \item $x^{2} \, u_{xx} + 2 \, x \, y \, u_{xy} + y^{2} \, u_{yy} = 0$
            \item $y^{2} \, u_{xx} + 2 \, x \, y \, u_{xy} + \left( x^{2} + 4 \, x^{4} \right) \, u_{yy} = 0$
        \end{enumerate}
        \item Considera la ecuación de calor:
            \begin{align*}
                u_{t} = \kappa \, u_{xx}
            \end{align*}
        sujeta a las CDF:
            \begin{align*}
            u(0, t) = 0 \hspace{1cm} u_{x} (L, t) = 0
            \end{align*}
        Resuelve el problema para las siguientes CI:
            \begin{enumerate}
                \item $u (x, 0) = 2 \, \cos \left( \dfrac{3 \pi x}{L} \right)$
                \item $u (x, 0) = 6 \sin \left( \dfrac{9 \pi x}{L} \right)$
                \item $u (x, 0) = 3 \sin \left( \dfrac{\pi x}{L} \right) - \sin \left( \dfrac{3 \pi x}{L} \right)$
                \item $u (x, 0) = \begin{cases}
                1, & \text{para } 0 < x < \dfrac{L}{2} \\[1em]
                2, & \text{para } \dfrac{L}{2} < x < L
                \end{cases}$
            \end{enumerate}
    \end{enumerate}

\subsection*{Método de Frobenius.}

\begin{enumerate}
    \item Determina los puntos singulares de las siguientes ED, clasifica cada punto singular en regular o irregular.
        \begin{enumerate}
            \item $x^{3} \, \sderivada{y} + 4 \, x^{2} \, \pderivada{y} + 3 \, y = 0$
            \item $x \, \sderivada{y} - (x + 3)^{-2} \, y = 0$
            \item $(x^{2} - 9)^{2} \, \sderivada{y} + (x + 3) \, \pderivada{y} + 2 \, y = 0$
            \item $\sderivada{y} - \dfrac{1}{x} \, \pderivada{y} + \dfrac{1}{(x - 1)^{3}} \, y = 0$
        \end{enumerate}
    \item Resuelve las siguientes ED con el método de Frobenius:
        \begin{enumerate} 
            \item $2 \, x \, \sderivada{y} - \pderivada{y} + 2 \, y = 0$
            \item $2 \, x \, \sderivada{y} + 5 \, \pderivada{y} + x \, y = 0$
            \item $x (x - 1) \, \sderivada{y} + 3 \, \pderivada{y} - 2 \, y = 0$
            \item $\sderivada{y} - \dfrac{3}{x} \, \pderivada{y} - 2 \, y = 0$
        \end{enumerate}
    \item Con la siguiente ecuación:
        \begin{align*}
            x^{2} \, \sderivada{y} (x) + x \, \pderivada{y} (x) + (x^{2} - 4) \, y (x) = 0
        \end{align*}
        Determina:
        \begin{enumerate}
            \item La naturaleza de los puntos singulares.
            \item La ecuación de índices.
            \item La regla de recurrencia.
            \item La(s) solución(es) en series de potencias.
        \end{enumerate}
\end{enumerate}

\subsection*{Segunda solución linealmente independiente.}

\begin{enumerate}
    \item Considerando que una solución de:
    \begin{align*}
        \sderivada{R} + \dfrac{1}{r} \, \pderivada{R} - \dfrac{m^{2}}{r^{2}} \, R = 0
    \end{align*}
    es $R = r^{m}$.
    
    Demuestra que la ecuación:
    \begin{align*}
        y_{2} (x) =  y_{1} \: (x) \scaleint{6ex}^{x} \dfrac{\exp \left[ \displaystyle - \scaleint{6ex}^{x_{2}} P (x_{1}) \: \dd{x_{1}} \right]}{[y_{1} (x_{2})]^{2}} \dd{x_{2}}
    % \label{eq:ecuacion_09_127}
    \end{align*}
    predice una segunda solución: $R = r^{-m}$
\end{enumerate}

\end{document}